\documentclass[a4paper,12pt]{article}
\usepackage[spanish]{babel}
\usepackage[utf8]{inputenc}
\usepackage{amsmath,amssymb,amsthm,mathtools}
\usepackage{geometry}
\geometry{margin=2.5cm}
\usepackage{enumitem}

\usepackage[dvipsnames]{xcolor}
\usepackage{tcolorbox}
\tcbuselibrary{skins,breakable}
\usepackage{titlesec}
\usepackage[hidelinks,colorlinks=true,linkcolor=NavyBlue,urlcolor=NavyBlue,citecolor=NavyBlue]{hyperref}
\usepackage{listings}
\lstset{language=Java, basicstyle=\ttfamily\small, keywordstyle=\color{NavyBlue}\bfseries,
  commentstyle=\color{ForestGreen}, stringstyle=\color{Maroon},
  showstringspaces=false, breaklines=true, frame=single, backgroundcolor=\color{gray!6},
  rulecolor=\color{gray!60}, tabsize=2}

\theoremstyle{definition}
\newtheorem{ejercicio}{Ejercicio}[section]

\tcolorboxenvironment{ejercicio}{%
  enhanced, breakable, colback=TealBlue!6, colframe=TealBlue!80!black,
  boxrule=1.1pt, arc=3mm, left=3mm, right=3mm, top=2mm, bottom=2mm,
  borderline west={3pt}{0pt}{TealBlue!90!black}%
}

\titleformat{\section}{\Large\bfseries\color{NavyBlue}}{\thesection}{1em}{}[\vspace{1mm}\color{NavyBlue}\titlerule]
\titleformat{\subsection}{\large\bfseries\color{MidnightBlue}}{\thesubsection}{1em}{}

\setlength{\parskip}{6pt}
\setlength{\parindent}{0pt}

\begin{document}

\begin{tcolorbox}[enhanced, colback=NavyBlue, colframe=NavyBlue!80!black,
  arc=4mm, boxrule=1pt, left=4mm, right=4mm, top=3mm, bottom=3mm, halign=center]
{\Large\bfseries\color{white} Práctica de Matrices 2MF Arroyo Seco}\\[1mm]
{\color{white!90!} Prof.\ Javier Martín --- con aplicaciones en Java}
\end{tcolorbox}


%=====================================================
\section{Tipos de matrices: fila, columna, rectangular, cuadrada}
%=====================================================
\begin{ejercicio}[Notas y vectores]
$F = \begin{pmatrix} 8 & 7 & 9 \end{pmatrix}$ son las notas de un estudiante.
\begin{enumerate}[label=\alph*)]
\item ¿Qué orden tiene? ¿Es fila, columna, rectangular o cuadrada?
\item En Java sería \texttt{int[] notas = \{8,7,9\}}. Escribí el \texttt{for} que la recorre e imprime.
\end{enumerate}
\end{ejercicio}

\begin{ejercicio}[Color RGB]
$C = \begin{pmatrix} 255 \\ 128 \\ 0 \end{pmatrix}$.
\begin{enumerate}[label=\alph*)]
\item ¿Orden y tipo?
\item ¿Qué color representa? ¿Qué pasaría si la guardás como fila en vez de columna en tu programa?
\end{enumerate}
\end{ejercicio}

\begin{ejercicio}[Dataset de películas]
\[
A = \begin{pmatrix}
5 & 7 & 3 & 9 \\
8 & 6 & 4 & 2 \\
1 & 0 & 5 & 7
\end{pmatrix}
\]
3 usuarios $\times$ 4 películas.
\begin{enumerate}[label=\alph*)]
\item Orden y tipo. ¿Qué es $a_{23}$ y $a_{31}$ en palabras?
\item En Java: \texttt{int[][] a = new int[3][4]}. ¿Cuántos elementos guarda? ¿Qué índices tiene la última casilla?
\end{enumerate}
\end{ejercicio}

\begin{ejercicio}[Red de computadoras]
\[
G = \begin{pmatrix}
0 & 1 & 1 \\
1 & 0 & 0 \\
1 & 0 & 0
\end{pmatrix}
\]
\begin{enumerate}[label=\alph*)]
\item ¿Cuadrada de qué orden? Escribí su diagonal principal.
\item ¿Qué pares de computadoras están conectados? Dibujá el grafo de 3 nodos.
\end{enumerate}
\end{ejercicio}

%=====================================================
\section{Nula, identidad y triangulares}
%=====================================================
\begin{ejercicio}[Imagen negra]
Una imagen $2 \times 3$ toda negra es $O_{2 \times 3}$.
\begin{enumerate}[label=\alph*)]
\item Escribila. ¿Qué orden tiene?
\item En Java, ¿cómo inicializás \texttt{int[][] img = new int[2][3]} en cero con dos \texttt{for}?
\end{enumerate}
\end{ejercicio}

\begin{ejercicio}[Identidad]
\begin{enumerate}[label=\alph*)]
\item Escribí $I_2$ e $I_3$.
\item En Java, construí $I_4$ con: $1$ si $i==j$, $0$ si no. Escribí el doble \texttt{for} con el \texttt{if}.
\end{enumerate}
\begin{lstlisting}
// pista:
for (int i = 0; i < 4; i++)
  for (int j = 0; j < 4; j++)
    // completar
\end{lstlisting}
\end{ejercicio}

\begin{ejercicio}[¿Triangular?]
\[
S = \begin{pmatrix} 2 & 5 & 1 \\ 0 & 3 & 4 \\ 0 & 0 & 6 \end{pmatrix}, \quad
L = \begin{pmatrix} 1 & 0 & 0 \\ 2 & 5 & 0 \\ 3 & -1 & 4 \end{pmatrix}, \quad
M = \begin{pmatrix} 1 & 2 & 0 \\ 0 & 1 & 0 \\ 1 & 0 & 1 \end{pmatrix}.
\]
Decí si cada una es triangular superior, inferior, diagonal o ninguna.
\end{ejercicio}

\begin{ejercicio}[Dependencias]
El módulo 2 depende del 1, y el 3 depende del 1 y 2. La matriz de dependencias es triangular inferior con $1$ en la diagonal.
\begin{enumerate}[label=\alph*)]
\item Escribí esa matriz $3 \times 3$.
\item ¿Qué significan los ceros sobre la diagonal en términos del programa?
\end{enumerate}
\end{ejercicio}

%=====================================================
\section{Traza}
%=====================================================
\begin{ejercicio}[Cálculo directo]
Calculá la traza de $A = \begin{pmatrix} 4 & 1 \\ 2 & 3 \end{pmatrix}$ y de $B = \begin{pmatrix} 1 & 0 & 2 \\ 5 & -1 & 3 \\ 0 & 4 & 6 \end{pmatrix}$.
\end{ejercicio}

\begin{ejercicio}[Matriz de confusión]
\[
C = \begin{pmatrix}
50 & 2 & 1 \\
3 & 45 & 5 \\
0 & 4 & 60
\end{pmatrix}
\]
\begin{enumerate}[label=\alph*)]
\item Calculá $\operatorname{tr}(C)$. ¿Qué mide (aciertos totales)?
\item En Java, ¿cómo sumarías la diagonal con un solo \texttt{for}?
\end{enumerate}
\begin{lstlisting}
// pista:
int traza = 0;
for (int i = 0; i < 3; i++) traza += c[i][i];
\end{lstlisting}
\end{ejercicio}

\begin{ejercicio}[Bucles en un grafo]
$G = \begin{pmatrix} 0 & 1 & 0 \\ 0 & 0 & 1 \\ 1 & 0 & 1 \end{pmatrix}$.
\begin{enumerate}[label=\alph*)]
\item Calculá $\operatorname{tr}(G)$. ¿Hay nodos conectados consigo mismos? ¿Cuáles?
\item ¿Qué significa $\operatorname{tr}(G)=0$ en una red?
\end{enumerate}
\end{ejercicio}

%=====================================================
\section{Igualdad, suma y producto por escalar}
%=====================================================
\begin{ejercicio}[Test de igualdad]
\begin{enumerate}[label=\alph*)]
\item ¿Son iguales $\begin{pmatrix} 1 & 2 \\ 3 & 4 \end{pmatrix}$ y $\begin{pmatrix} 1 & 2 \\ 3 & 4 \end{pmatrix}$? ¿Y con $\begin{pmatrix} 1 & 2 & 0 \\ 3 & 4 & 0 \end{pmatrix}$?
\item Hallá $x,y$: $\begin{pmatrix} x & 2 \\ 3 & y \end{pmatrix} = \begin{pmatrix} 1 & 2 \\ 3 & 5 \end{pmatrix}$.
\item En Java, ¿cómo compararías dos matrices $2 \times 2$ casilla por casilla? Escribí la idea con \texttt{boolean iguales}.
\end{enumerate}
\end{ejercicio}

\begin{ejercicio}[Turnos de una tienda]
Mañana $M = \begin{pmatrix} 10 & 5 & 3 \\ 4 & 7 & 2 \end{pmatrix}$, tarde $T = \begin{pmatrix} 6 & 1 & 8 \\ 3 & 9 & 5 \end{pmatrix}$ (2 días $\times$ 3 productos). Calculá el total $M+T$ y explicá qué es cada casilla.
\end{ejercicio}

\begin{ejercicio}[Evento puntaje triple]
Recompensas $R = \begin{pmatrix} 10 & 25 & 30 \\ 15 & 5 & 20 \end{pmatrix}$. Hay evento $\times 3$.
\begin{enumerate}[label=\alph*)]
\item Calculá $3R$.
\item Normalizá la imagen $\begin{pmatrix} 255 & 128 \\ 0 & 64 \end{pmatrix}$ con $\frac{1}{255}A$.
\item En Java, multiplicá toda matriz por \texttt{k} sin crear otra matriz (en el lugar):
\end{enumerate}
\begin{lstlisting}
for (int i = 0; i < f; i++)
  for (int j = 0; j < c; j++)
    m[i][j] = m[i][j] * k;
\end{lstlisting}
\end{ejercicio}

%=====================================================
\section{Fila por columna y producto de matrices}
%=====================================================
\begin{ejercicio}[Promedio ponderado]
Notas $\begin{pmatrix} 7 & 8 & 9 \end{pmatrix}$, pesos $\begin{pmatrix} 0{,}3 \\ 0{,}3 \\ 0{,}4 \end{pmatrix}$.
\begin{enumerate}[label=\alph*)]
\item Calculá el producto. ¿Qué número da?
\item En Java es el producto punto. Completalo:
\end{enumerate}
\begin{lstlisting}
int suma = 0;
for (int k = 0; k < 3; k++) suma += fila[k] * col[k];
\end{lstlisting}
\end{ejercicio}

\begin{ejercicio}[Neurona]
$w = \begin{pmatrix} 2 & -1 & 4 \end{pmatrix}$, $x = \begin{pmatrix} 3 \\ 0 \\ -2 \end{pmatrix}$. Calculá $w \cdot x$ (es $y = wx$ sin sesgo).
\end{ejercicio}

\begin{ejercicio}[Producto paso a paso]
\[
A = \begin{pmatrix} 1 & 2 & 3 \\ 4 & 5 & 6 \end{pmatrix}_{2 \times 3}, \quad
B = \begin{pmatrix} 7 & 8 \\ 9 & 10 \\ 11 & 12 \end{pmatrix}_{3 \times 2}.
\]
\begin{enumerate}[label=\alph*)]
\item ¿Se puede hacer $AB$? ¿De qué orden es? Calculalo completo.
\item ¿Se puede hacer $BA$? Si sí, ¿de qué orden?
\item Explicá por qué $A_{2 \times 3}$ por $D_{2 \times 2}$ es imposible (como el error \texttt{shapes not aligned}).
\end{enumerate}
\end{ejercicio}

\begin{ejercicio}[Capa lineal]
$W = \begin{pmatrix} 2 & -1 \\ 0 & 3 \end{pmatrix}$, $x = \begin{pmatrix} 4 \\ 5 \end{pmatrix}$. Calculá $Wx$ y verificá que da $\begin{pmatrix} 3 \\ 15 \end{pmatrix}$.
\end{ejercicio}

%=====================================================
\section{Propiedades, potencia e idempotencia}
%=====================================================
\begin{ejercicio}[Asociativa y distributiva]
Con $A=\begin{pmatrix} 1 & 2 \\ 0 & 1 \end{pmatrix}$, $B=\begin{pmatrix} 2 & 0 \\ 1 & 1 \end{pmatrix}$, $C=\begin{pmatrix} 1 \\ 3 \end{pmatrix}$ verificá $(AB)C = A(BC)$. Luego, con otras $2 \times 2$ a elección, verificá $A(B+C)=AB+AC$.
\end{ejercicio}

\begin{ejercicio}[El orden importa]
$R=\begin{pmatrix} 0 & -1 \\ 1 & 0 \end{pmatrix}$ (rota 90$^\circ$), $E=\begin{pmatrix} 2 & 0 \\ 0 & 1 \end{pmatrix}$ (estira en $x$).
\begin{enumerate}[label=\alph*)]
\item Calculá $RE$ y $ER$. ¿Son iguales?
\item Aplicalas al punto $\begin{pmatrix} 1 \\ 1 \end{pmatrix}$. ¿Qué puntos distintos obtenés?
\item En un juego, ¿por qué el orden del \textit{pipeline} cambia el dibujo?
\end{enumerate}
\end{ejercicio}

\begin{ejercicio}[Potencias]
$A=\begin{pmatrix} 1 & 1 \\ 1 & 0 \end{pmatrix}$. Calculá $A^2$ y $A^3$. Para $D=\begin{pmatrix} 2 & 0 \\ 0 & 3 \end{pmatrix}$ calculá $D^2$ y $D^3$ (pista: en diagonales se potencia directo).
\end{ejercicio}

\begin{ejercicio}[Filtro estable e idempotente]
$P=\begin{pmatrix} 1 & 0 \\ 0 & 0 \end{pmatrix}$ y $Q=\begin{pmatrix} 0{,}5 & 0{,}5 \\ 0{,}5 & 0{,}5 \end{pmatrix}$.
\begin{enumerate}[label=\alph*)]
\item Verificá $P^2=P$ y $Q^2=Q$.
\item Explicá con palabras: ¿por qué aplicar dos veces el filtro es igual que una?
\end{enumerate}
\end{ejercicio}

%=====================================================
\section{Regular, singular, rango e inversa por Gauss-Jordan}
%=====================================================
\begin{ejercicio}[¿Regular o singular?]
$A=\begin{pmatrix} 1 & 2 \\ 3 & 4 \end{pmatrix}$, $B=\begin{pmatrix} 1 & 2 \\ 2 & 4 \end{pmatrix}$.
\begin{enumerate}[label=\alph*)]
\item Sin calcular la inversa, ¿cuál sospechás singular (fila múltiplo de otra)?
\item Verificá que $A^{-1}=\begin{pmatrix} -2 & 1 \\ 1{,}5 & -0{,}5 \end{pmatrix}$ cumple $AA^{-1}=I_2$.
\end{enumerate}
\end{ejercicio}

\begin{ejercicio}[Inversa $2 \times 2$]
Por Gauss-Jordan, hallá la inversa de $\begin{pmatrix} 1 & 2 \\ 3 & 4 \end{pmatrix}$ partiendo de $(A\mid I_2)$. Mostrá las 3 operaciones ($F_2-3F_1$, dividir, $F_1-2F_2$).
\end{ejercicio}

\begin{ejercicio}[Inversa $3 \times 3$]
Hallá la inversa de $A=\begin{pmatrix} 1 & 0 & 1 \\ 0 & 1 & 1 \\ 1 & 1 & 0 \end{pmatrix}$ con $(A\mid I_3)$. Respuesta: $\begin{pmatrix} 1/2 & -1/2 & 1/2 \\ -1/2 & 1/2 & 1/2 \\ 1/2 & 1/2 & -1/2 \end{pmatrix}$. Verificala con $AA^{-1}=I_3$.
\end{ejercicio}

\begin{ejercicio}[Inversa $3 \times 3$ triangular]
Por Gauss-Jordan, hallá la inversa de $C=\begin{pmatrix} 2 & 0 & 0 \\ 1 & 3 & 0 \\ 0 & 2 & 1 \end{pmatrix}$ partiendo de $(C\mid I_3)$.
\end{ejercicio}

\begin{ejercicio}[Sin inversa: $3 \times 3$ singular]
Intentá hallar por Gauss-Jordan la inversa de $B=\begin{pmatrix} 1 & 2 & 3 \\ 0 & 1 & 1 \\ 1 & 3 & 4 \end{pmatrix}$ con $(B\mid I_3)$.
\begin{enumerate}[label=\alph*)]
\item Hacé $F_3 \to F_3-F_1-F_2$. ¿Qué fila obtenés a la izquierda?
\item ¿Por qué eso prueba que $B$ es singular y no tiene inversa? ¿Qué rango tiene $B$?
\item En Java, si tu función \texttt{inversa()} llega a una fila $0\cdots 0$ a la izquierda, ¿qué debería devolver o informar?
\end{enumerate}
\end{ejercicio}

\begin{ejercicio}[Otra singular: filas proporcionales]
$M=\begin{pmatrix} 1 & 0 & 2 \\ 2 & 1 & 1 \\ 3 & 1 & 3 \end{pmatrix}$ ($F_3=F_1+F_2$).
\begin{enumerate}[label=\alph*)]
\item Escaloná $(M\mid I_3)$ hasta obtener una fila nula a la izquierda. Concluí que no existe $M^{-1}$.
\item ¿Qué relación hay entre ``filas dependientes'' y ``columnas redundantes en un dataset''?
\end{enumerate}
\end{ejercicio}

\begin{ejercicio}[Rango = filas útiles]
Escaloná $\begin{pmatrix} 1 & 2 \\ 2 & 4 \end{pmatrix}$ con $F_2-2F_1$. ¿Cuántos escalones quedan? ¿Rango? Explicá por qué en Java eso sería un sensor redundante (dos filas que miden lo mismo).
\end{ejercicio}

%=====================================================
\section{Traspuesta y matrices especiales}
%=====================================================
\begin{ejercicio}[Trasponer dataset]
$H=\begin{pmatrix} 5 & 6 & 7 \\ 8 & 9 & 10 \end{pmatrix}_{2 \times 3}$ (2 semanas $\times$ 3 cursos).
\begin{enumerate}[label=\alph*)]
\item Calculá $H^t$. ¿Qué orden tiene?
\item En Java completalo (fijate que las dimensiones se invierten):
\end{enumerate}
\begin{lstlisting}
int[][] t = new int[3][2];
for (int i = 0; i < 2; i++)
  for (int j = 0; j < 3; j++)
    t[j][i] = h[i][j];
\end{lstlisting}
\end{ejercicio}

\begin{ejercicio}[Las 4 propiedades]
Con $A=\begin{pmatrix} 1 & 2 \\ 3 & 4 \end{pmatrix}$, $B=\begin{pmatrix} 0 & 1 \\ -1 & 2 \end{pmatrix}$ verificá $(A+B)^t=A^t+B^t$ y $(AB)^t=B^tA^t$. Con $k=2$ verificá $(kA)^t=kA^t$ y $(A^t)^t=A$.
\end{ejercicio}

\begin{ejercicio}[Clasificar]
\[
S=\begin{pmatrix} 0 & 1 & 1 \\ 1 & 0 & 1 \\ 1 & 1 & 0 \end{pmatrix}, \;
K=\begin{pmatrix} 0 & 2 & -1 \\ -2 & 0 & 3 \\ 1 & -3 & 0 \end{pmatrix}, \;
R=\begin{pmatrix} 0 & -1 \\ 1 & 0 \end{pmatrix}.
\]
\begin{enumerate}[label=\alph*)]
\item ¿Cuál es simétrica ($A^t=A$), antisimétrica ($A^t=-A$) y ortogonal ($R^tR=I$)?
\item Asociá cada una a: red no dirigida, torneo con $a_{ij}=-a_{ji}$, rotación de videojuego.
\item En Java, ¿cómo chequearías simetría con \texttt{m[i][j]==m[j][i]} en doble \texttt{for}?
\end{enumerate}
\end{ejercicio}

%=====================================================
\section{Rouché-Frobenius y ecuaciones matriciales}
%=====================================================
\begin{ejercicio}[Tres sistemas $3 \times 3$, tres destinos]
Escaloná $(A\mid B)$ y clasificá con Rouché-Frobenius. Todos son $3 \times 3$ (3 ecuaciones, 3 incógnitas):
\begin{enumerate}[label=\alph*)]
\item Determinado:
$\left(\begin{array}{ccc|c} 1 & 0 & 1 & 3 \\ 0 & 1 & 1 & 4 \\ 1 & 1 & 0 & 3 \end{array}\right)$
(pista: $F_3 \to F_3-F_1$, luego $F_3 \to F_3-F_2$; quedan 3 escalones).
\item Incompatible:
$\left(\begin{array}{ccc|c} 1 & 2 & 3 & 1 \\ 0 & 1 & 1 & 2 \\ 1 & 3 & 4 & 5 \end{array}\right)$
(pista: $F_3 \to F_3-F_1-F_2$ da la fila $0\;0\;0\mid 2$, o sea $0=2$).
\item Indeterminado:
$\left(\begin{array}{ccc|c} 1 & 2 & 3 & 1 \\ 0 & 1 & 1 & 2 \\ 1 & 3 & 4 & 3 \end{array}\right)$
(pista: la misma operación da $0\;0\;0\mid 0$; la tercera ecuación sobra).
\end{enumerate}
Para cada uno da $\operatorname{rango}(A)$, $\operatorname{rango}(A\mid B)$ y si es compatible determinado / incompatible / indeterminado. En b) y c) la matriz $A$ es la misma singular de rango 2: lo que cambia el destino es $B$.
\end{ejercicio}

\begin{ejercicio}[Resolver $AX=B$]
$A=\begin{pmatrix} 1 & 2 \\ 3 & 4 \end{pmatrix}$, $B=\begin{pmatrix} 5 \\ 11 \end{pmatrix}$. Con $A^{-1}=\begin{pmatrix} -2 & 1 \\ 1{,}5 & -0{,}5 \end{pmatrix}$ hallá $X=A^{-1}B$ y comprobá que es $\begin{pmatrix} 1 \\ 2 \end{pmatrix}$.
\end{ejercicio}

\begin{ejercicio}[Decodificar $XA=B$]
$A=\begin{pmatrix} 2 & 0 \\ 0 & 3 \end{pmatrix}$ escaló por 2 y 3, $B=\begin{pmatrix} 8 & 12 \\ 4 & 6 \end{pmatrix}$. Hallá $X=BA^{-1}$ con $A^{-1}=\begin{pmatrix} 1/2 & 0 \\ 0 & 1/3 \end{pmatrix}$. Interpretá: $X$ es la imagen original antes del escalado.
\end{ejercicio}

\begin{ejercicio}[Desafío $AXB=C$]
$A=\begin{pmatrix} 1 & 1 \\ 0 & 1 \end{pmatrix}$, $B=\begin{pmatrix} 1 & 0 \\ 1 & 1 \end{pmatrix}$, $C=\begin{pmatrix} 3 & 2 \\ 1 & 1 \end{pmatrix}$. Con $A^{-1}=\begin{pmatrix} 1 & -1 \\ 0 & 1 \end{pmatrix}$, $B^{-1}=\begin{pmatrix} 1 & 0 \\ -1 & 1 \end{pmatrix}$ hallá $X=A^{-1}CB^{-1}$ y verificá $AXB=C$.
\end{ejercicio}

%=====================================================
\section*{Respuestas (cortas)}
%=====================================================
\begin{enumerate}[label=\textbf{\arabic*.}, leftmargin=*]
\item[1.] 1.1: $1\times 3$, fila. 1.2: $3\times 1$, columna, naranja. 1.3: $3\times 4$ rectangular, $a_{23}=4$, $a_{31}=1$, 12 elementos, última $(2,3)$ en Java. 1.4: orden 3, diagonal $0,0,0$, conexiones 1--2 y 1--3.
\item[2.] 2.1: $O_{2\times 3}$ todo cero. 2.2: $I_2,I_3$; código con \texttt{if (i==j) 1 else 0}. 2.3: $S$ sup., $L$ inf., $M$ ninguna. 2.4: $\begin{pmatrix}1&0&0\\1&1&0\\1&1&1\end{pmatrix}$; ceros = sin dependencia hacia adelante.
\item[3.] 3.1: $7$ y $6$. 3.2: $155$ aciertos. 3.3: $\operatorname{tr}=1$, bucle en nodo 3; $\operatorname{tr}=0$ = sin bucles.
\item[4.] 4.1: sí; no (distinto orden); $x=1,y=5$. 4.2: $\begin{pmatrix}16&6&11\\7&16&7\end{pmatrix}$. 4.3: $3R=\begin{pmatrix}30&75&90\\45&15&60\end{pmatrix}$; normalizada $\begin{pmatrix}1&0{,}50\\0&0{,}25\end{pmatrix}$.
\item[5.] 5.1: $8{,}1$. 5.2: $-2$. 5.3: $AB=\begin{pmatrix}58&64\\139&154\end{pmatrix}_{2\times2}$; $BA$ $3\times3$ sí existe; $2\times3$ por $2\times2$ imposible. 5.4: $\begin{pmatrix}3\\15\end{pmatrix}$.
\item[6.] 6.1: $(AB)C=A(BC)=\begin{pmatrix}13\\4\end{pmatrix}$. 6.2: $RE=\begin{pmatrix}0&-1\\2&0\end{pmatrix}\neq ER=\begin{pmatrix}0&-2\\1&0\end{pmatrix}$; puntos $(-1,2)$ vs $(-2,1)$. 6.3: $A^2=\begin{pmatrix}2&1\\1&1\end{pmatrix}$, $A^3=\begin{pmatrix}3&2\\2&1\end{pmatrix}$; $D^2=\begin{pmatrix}4&0\\0&9\end{pmatrix}$. 6.4: $P^2=P$, $Q^2=Q$.
\item[7.] 7.1: $A$ regular, $B$ singular. 7.2--7.3: ver teoremas del teórico. Triangular $C^{-1}=\begin{pmatrix}1/2&0&0\\-1/6&1/3&0\\1/3&-2/3&1\end{pmatrix}$. $B$ y $M$ singulares: aparece fila $0\;0\;0$ a la izquierda, rango 2, sin inversa. 7.4: rango 1.
\item[8.] 8.1: $H^t_{3\times2}$. 8.2: verificación directa. 8.3: $S$ simétrica (red), $K$ antisimétrica (torneo), $R$ ortogonal (rotación).
\item[9.] 9.1: a) rangos $3,3=n$ determinado, solución $x=1,y=2,z=2$; b) $2\neq3$ incompatible ($0=2$); c) $2=2<3$ indeterminado (infinitas, 1 grado de libertad). 9.2: $\begin{pmatrix}1\\2\end{pmatrix}$. 9.3: $\begin{pmatrix}4&4\\2&2\end{pmatrix}$. 9.4: $\begin{pmatrix}1&1\\0&0\end{pmatrix}$.
\end{enumerate}

\end{document}
